% No se simulan negritas ni cursivas: las cuatro variantes deben existir.
\defaultfontfeatures{Ligatures=TeX}
\newcommand{\SRLLoadSystemFont}{%
  \IfFontExistsTF{\SRLFontName}{%
    \setmainfont{\SRLFontName}%
    \edef\SRLActualFont{\SRLFontName}%
    \typeout{SRL-FONT-SOURCE: system}%
  }{%
    \setmainfont{\SRLFallbackFont}%
    \edef\SRLActualFont{\SRLFallbackFont}%
    \PackageWarningNoLine{srl-template}{Fuente solicitada '\SRLFontName' no disponible.
      Se usa '\SRLFallbackFont'. Esto NO es una compilacion con Minion Pro}%
    \typeout{SRL-FONT-SOURCE: fallback}%
  }%
}
\newcommand{\SRLLoadLocalFont}{%
  \setmainfont{\SRLFontRegular}[
    Path=\SRLFontDirectory,
    UprightFont=\SRLFontRegular,
    ItalicFont=\SRLFontItalic,
    BoldFont=\SRLFontBold,
    BoldItalicFont=\SRLFontBoldItalic
  ]%
  \edef\SRLActualFont{\SRLFontName}%
  \typeout{SRL-FONT-SOURCE: files}%
}
\ifdefstrequal{\SRLFontName}{\SRLLocalFontFamily}{%
  \IfFileExists{\SRLFontDirectory\SRLFontRegular}
    {\SRLLoadLocalFont}{\SRLLoadSystemFont}%
}{\SRLLoadSystemFont}
\AtBeginDocument{\typeout{SRL-FONT: \SRLActualFont}}
